% What is the Gauge RR? 
\begin{frame}
    [t]
    {What is the Gauge RR approach?}
    \fst{Concept and Problems of Classical Approach}
    \vspace{0.5em}
    \tv{Goal of Gauge RR Study}\\
    \begin{itemize}
        \footnotesize
        \item To understand and quantify the sources variability\\
        {\scriptsize $\rightarrow$ \tv{Estimates of variation, identify the most influential variation factors}\par}
        \item To improve and control measurement process
        \item To compensate the variability in evaluating the product
    \end{itemize}
    \vspace{1em}
    \tv{Classical Model of Gauge RR}
    \vspace{-1em}
    {\footnotesize \begin{equation*}
        X_{ijk} = \mu + O_i + P_j + {(OP)}_{ij} + R_{k(ij)}\qquad
        \begin{cases}
            i = 1, 2, \ldots, o;\\
            j = 1, 2, \ldots, p;\\
            k = 1, 2, \ldots, n;
        \end{cases}
    \end{equation*}}
    \vspace{-1em}
    \begin{itemize}
        \footnotesize
        \item Parts and Operator should be considered as '\tv{random factor}'.
        \item Repeatablity factor ($R$) is \tv{nested} in Parts ($P$) and Operator ($O$) factor. 
    \end{itemize}

\end{frame}